% One-time imported draft. Editable native source; no Markdown regeneration.
\section{Express the rule as a Boolean
calculation}\label{008-life-generations-p07-express-the-rule-as-a-boolean-calculation}

A next state is alive when the neighbor count is three, or when the
count is two and the current cell is alive. The second clause matters:
two neighbors do not create a new cell from a dead center.

The reference turns equality into zero with XOR, then turns zero into
one with an unsigned less-than-one test. Those tests yield binary flags.
AND gates the two-neighbor flag with the current alive byte, and OR
combines that result with the three-neighbor flag. The final value is
exactly zero or one.

Branches can express the same truth table more directly. The alternative
chooses zero by default, chooses one for a count of three, and copies
the current state for a count of two. Instruction structure can vary
while the local rule and buffer-phase contract remain unchanged. The
rule table is the common meaning behind either source.
